% Frugal AI seminar, part 1: why frugal AI. Every figure is one a cited source states (named in the Ref line).
% Sources are listed in ai_frugal_join.tex. CheatSheet-safe: no floats.

\section[Why]{Why Frugal AI?}

%%%%%%%%%%%%%%%%%%%%%%%%%%%%%%%%%%%%%%%%%%%%%%%%%%%
\begin{frame}[fragile]\frametitle{}
\begin{center}
{\Large Why Frugal AI?}

\medskip
{\small (The least that does the job)}
\end{center}
\end{frame}

%%%%%%%%%%%%%%%%%%%%%%%%%%%%%%%%%%%%%%%%%%%%%%%%%%%
\begin{frame}[fragile]\frametitle{The Default Answer: Bigger}
\begin{itemize}
\item The usual path today: pick the biggest model, call it through a hosted API, pay per token.
\item That works, until the bill, the delay, the connection or the data rules get in the way.
\item Frugal AI asks the opposite question: what is the smallest thing that does this job?
\item It is not a worse AI. It is a fitted one: the right size, in the right place, for the task.
\item This talk: why it matters, how small models work, where they run, what already works, and how to start.
\end{itemize}
\end{frame}

%%%%%%%%%%%%%%%%%%%%%%%%%%%%%%%%%%%%%%%%%%%%%%%%%%%
\begin{frame}[fragile]\frametitle{What One Prompt Costs}
\begin{center}
\adjustbox{max width=0.8\linewidth, max totalheight=0.40\textheight}{\includegraphics{frugal_energy_per_prompt}}
\end{center}
\begin{itemize}
\item Google measured the median Gemini Apps text prompt at 0.24 Wh, counting host machines, idle capacity and data centre overhead: its own service, its own boundary.
\item A benchmark of 30 commercial models estimates over 29 Wh for a long prompt on the most energy-intensive ones, from public API data: an estimate, not a measurement.
\item Different methods and different prompts: read the range, not a ranking.
\end{itemize}
{\tiny (Ref: Elsworth et al., arXiv 2508.15734; Jegham et al., arXiv 2505.09598)}
\end{frame}

%%%%%%%%%%%%%%%%%%%%%%%%%%%%%%%%%%%%%%%%%%%%%%%%%%%
\begin{frame}[fragile]\frametitle{Data Centres in Context}
\begin{itemize}
\item The IEA reports that data centres used about 415 TWh in 2024, around 1.5\% of world electricity.
\item It projects about 945 TWh by 2030, more than double.
\item These are projections of the whole sector, not a measurement of any one model or prompt.
\item So the lesson is not ``AI is bad'': it is that every token has a cost, and the cheapest token is the one you did not need.
\end{itemize}
{\tiny (Ref: IEA, Energy and AI, 2025, executive summary)}
\end{frame}

%%%%%%%%%%%%%%%%%%%%%%%%%%%%%%%%%%%%%%%%%%%%%%%%%%%
\begin{frame}[fragile]\frametitle{What Is Frugal AI?}
\begin{itemize}
\item The Frugal AI Hub defines it as an approach that ``maximizes environmental and societal impact with minimal resources''.
\item Its own example: a small, locally hosted language model on a Raspberry Pi, giving offline healthcare in rural clinics.
\item A July 2026 report from the Hub describes India's strength as ``AI that works cheaply, offline and in-language at population scale''.
\item Four words to remember: efficient, affordable, sustainable, accessible.
\end{itemize}
{\tiny (Ref: Frugal AI Hub glossary; Frugal AI Hub, Minds, Models \& Markets, July 2026)}
\end{frame}

%%%%%%%%%%%%%%%%%%%%%%%%%%%%%%%%%%%%%%%%%%%%%%%%%%%
\begin{frame}[fragile]\frametitle{Three Levers}
\begin{itemize}
\item \textbf{Smaller model}: pick the size the job needs, then shrink it further (quantisation, distillation).
\item \textbf{Closer hardware}: run on the laptop, board or phone you already own, offline if the job allows.
\item \textbf{Fewer tokens}: lean context, short structured prompts, step budgets, no re-sending of whole histories.
\item A study of task-specific systems against multi-purpose generative models finds the multi-purpose ones orders of magnitude more expensive per 1,000 inferences on its benchmark tasks. That is its result, not a law for every job.
\end{itemize}
{\tiny (Ref: Luccioni, Jernite, Strubell, Power Hungry Processing, FAccT 2024)}
\end{frame}
